\cotefile{136}
\scanpage{64}

% top, opening two lines
Mais si $F$ \inserted{sur $X_*$} est \deleted{loc. constant}\inserted{un $k$-Module} \inserted{($k=$ un anneau \unsure{constant comm.|commutatif})}, alors on vérifie facilement que

% upper, first display and its underbrace
\[
\unsure{C^{*}|C^{\bullet}}(X_*,F)\simeq\Gamma\bigl(X_*,\underbrace{(C^{*}_{*}\mid\unsure{X_{*}|XX_{*}})}_{C^{*}_{X_*}}\otimes_{k_*}F\bigr)
\]

% left, encircled beside first display
\margin{NB On pose $C^{*}_{*}=C^{*}_{*k}$ pour abréger}

% upper middle, resolution, marginal query and conclusion
et on vérifie tout de suite que $C^{*}_{X_\bullet}\otimes_{k_\bullet}F$ est une résolution de $F$, laquelle est d'ailleurs \uline{flasque}\margin{à vérifier ce point $\longrightarrow$} \deleted{Donc} (que $F$ soit loc. constant ou non, d'ailleurs). Donc, sans hypothèse de locale constance sur $F$, $C^{*}(X_\bullet,F)$ permet de calculer la coh. de $X_\bullet$ à coéff. dans $F$ (i.e. on a un isom. can.

% middle, cohomology display and definition beneath C
\[
H^{*}\bigl(\underbrace{C^{*}(X_*,F]}_{\substack{\Vert\,\mathrm{dfn}\\\Gamma(X_\bullet,C^{*}_{X_\bullet}\otimes_k F}}\bigr)\simeq H^{*}(X_\bullet,F)
\]

% middle to lower, functoriality and hypercohomology
pour \deleted{\unsure{compatible}} fonctoriel en $F$ (et $=$ en $(X_\bullet,F)\ldots$) et compatible avec les hom. cobords. De $=$ pour le calcul de l'hypercohomologie de $X_\bullet$ à coeff. dans un complexe de $k$-Modules $F^{*}$. En termes de catégories dérivées, on aura un isomorphisme de foncteurs triangulés

% bottom, derived functor and explanation below brace
\[
R\Gamma(X_*,F^*)\simeq\underbrace{C^{*}(X_*,F^*)}_{\substack{\text{pour degré total}\\\text{associé aux}\\\text{deux degrés}\\\text{(\unsure{prov.|proven.|venant} de \unsure{\v{C} et|\v{C}ech|C\textsuperscript{*} et} de}\\ F^{*}\text{)}}}
\]

% lower left
% archive sheet number: 63
\scanpage{65}

% top, lines 1–3
Soit maintenant $A_*^{\unsure{0|\circ|*}}$ un \deleted{$k$-module} complexe de cochaines de $k$-modules $\frac{1}{2}$ simpliciaux

% upper, complex of A's
\[
0\mathrel{\unsure{\relbar|\to}} A_*^0\relbar A_\bullet^1\relbar A_\bullet^2\longrightarrow\ldots
\]

% upper, lines below first complex
et proposons-nous de déterminer les homomorphismes de ce complexe dans

% upper middle, complex of C's
\[
0\relbar C^0_{*M}\relbar C^1_{*M}\relbar C^2_{*M}\relbar\ldots,
\]

% middle, M and degree-zero maps
où $M$ est un $k$-module donné

Tout d'abord, les hom de degré $0$ (\unsure{pas|qui} \unsure{néc.} \unsure{commutant} avec op. diff) sont les syst. d'hom

% middle, system u
\[
A_*^n\xrightarrow{u_n}C^n_{*M}
\]

% middle, equivalent system
i.e. les syst. d'hom de $k$-modules

% middle, diagonal-degree system
\[
A_n^n\xrightarrow{\unsure{K_n|k_n|\kappa_n}}M
\]

% below middle, cancellation above L'hom
\deleted{\unsure{On a un|On a}} L'hom \unsure{composé|comp.}

% below middle, three horizontal objects and curved composite
\diagram{\begin{tikzcd}[ampersand replacement=\&,column sep=large]
d_C\circ u_n\mathpunct{:}\ A_\bullet^n\arrow[r,dash]\arrow[rr,bend right=16] \& C^n_{*M}\arrow[r] \& C^{n+1}_{*M}
\end{tikzcd}}

% lower, correspondence
correspond à \deleted{\unsure{un}} l'hom composé

% lower, first triangle
\diagram{\begin{tikzcd}[ampersand replacement=\&,column sep=small]
A^n_{n+1}\arrow[rr,"K_n\partial_{n+1}"]\arrow[dr,"\partial_{n+1}"'] \& \& M\\
\& A^n_n\arrow[ur,"K_n"'] \& {}
\end{tikzcd}}

% lower, comparison
tandis que l'hom composé

% lower, other composite
\[
u_{n+1}\circ d_A\mathpunct{:}\ A^n_*\longrightarrow A^{n+1}_*\longrightarrow C^{n+1}_{*M}
\]

% foot, lead to second triangle
correspond à l'hom composé

% foot, second triangle along the paper edge
\diagram{\begin{tikzcd}[ampersand replacement=\&,column sep=small]
A^n_{n+1}\arrow[rr]\arrow[dr,"d_A^n"'] \& \& M\\
\& A^{n+1}_{n+1}\arrow[ur,"K_{n+1}"'] \& {}
\end{tikzcd}}

% lower left
% archive sheet number: 64
\scanpage{66}

% top, lines 1–3
et la condition que les $u_n$ définissent bien un hom de complexes est la ``condition de Stokes''

% upper, Stokes relation beside its square
\diagram{$\unsure{K_n|k_n|\kappa_n}\partial_{n+1}=K_{n+1}d:\quad$
\begin{tikzcd}[ampersand replacement=\&]
A^n_{n+1}\arrow[r,"\partial_{n+1}"]\arrow[d,"d^n"'] \& A^n_n\arrow[d,"K_n"]\\
A^{n+1}_{n+1}\arrow[r,"K_{n+1}"'] \& M
\end{tikzcd}}

% upper middle, below square
Sous ces conditions, pour tt $X_*$\inserted{$\in\operatorname{Ob}\widehat{\Delta}$}, on a un hom. \deleted{\unsure{ca}} déduit de $u$

% middle, homomorphism and vertically identified target
\diagram{\begin{tikzcd}[ampersand replacement=\&,column sep=small,row sep=tiny]
\Gamma(X_*,A^*_{X_*})\arrow[r] \& C^{\unsure{*|\bullet}}(X_*,M_*)\\[-1ex]
\& \arrow[u,equal]\deleted{\unsure{C}}\Gamma(X_*,C^{*}_{\deleted{\unsure{X}}*M\deleted{\unsure{*}}}\mid X_\bullet)
\end{tikzcd}}

% middle, derived category
qui représente l'hom dans la cat. dérivée

% middle lower, derived-category map
\[
\Gamma(X_*,A^*_{X_*})\mathrel{\unsure{\relbar|\longrightarrow}}\mathbb{R}\Gamma_{X_\bullet}(M)
\]

% lower, before second square
associé à $u_*$. On aura alors commutativité dans

% lower, canonical square with cancelled first letter
\diagram{\begin{tikzcd}[ampersand replacement=\&,column sep=small]
\deleted{\mathbb{R}}\Gamma(X_\bullet,A^*_{X_*})\arrow[r]\arrow[d] \& C^*(X_\bullet,M)=\Gamma(X_\bullet,C^*_{M}\mid\unsure{X_\bullet|X_*})\arrow[d,"\sim"]\\
\mathbb{R}\Gamma(X_\bullet,A^*_{X_\bullet})\arrow[r] \& \mathbb{R}\Gamma(X_*,M)
\end{tikzcd}}

% foot, conditional sentence continues
Si $A^{*}$ est une résolution de $M_*$, et $K_0:A_0^0\longrightarrow M$ induit l'identité sur $M\subset A^0_0$, alors la \unsure{dernière|deuxième} flèche horizontale est aussi un iso, et

% lower left
% archive sheet number: 65
\scanpage{67}

% top, continuation
on peut considérer la flèche

% top, arrow with Hom identification below source
\diagram{\begin{tikzcd}[ampersand replacement=\&,row sep=tiny]
\Gamma(X_*,A^*_{X_\bullet})\arrow[r] \& C(X_*,M)\\
\operatorname{Hom}(X_\bullet,A^*)\arrow[u,equal] \& {}
\end{tikzcd}}

% upper, Stokes–De Rham system
déduite du syst. de Stokes-De Rham $(\unsure{K_n|k_n|\kappa_n})$ comme une réalisation de l'hom. canonique

% upper middle, derived global sections and upward arrow
\diagram{\begin{tikzcd}[ampersand replacement=\&,row sep=small]
\deleted{\unsure{R}}\Gamma(X_\bullet,A^*_{X_\bullet})\arrow[r] \& \mathbb{R}\Gamma(X_\bullet,A^*_{X_\bullet})\\
\& \mathbb{R}\Gamma(X_\bullet,M)\arrow[u,"\sim"']
\end{tikzcd}}

% middle, cancelled beginnings
\deleted{On note que sous les conditions de la prop. 1, il \unsure{y|a}} \deleted{Supposons}

% middle, NB
\uline{NB} L'hom

% middle, NB display with Hom at lower left
\[
\underset{\operatorname{Hom}(X_\bullet,A^n_*)}{\Gamma(X_*,A^n_{X_\bullet})}\xrightarrow{u^n} C^n(X_\bullet,M)
\]

% middle, explanation of component
associé à $K_n$ s'explicite ainsi : On considère, pour tt élément \deleted{\unsure{$u$ de}} $f:X_*\longrightarrow A^n_\bullet$, l'application

% below middle, component of f
\[
f_n:X_n\longrightarrow A^n_n
\]

% below middle, line above equation
et on prend

% lower, formula above separator
\[
u^n(f)=K_n\circ f_n\in\operatorname{Appl}(X_n,M)
\]

% lower, assumptions introduced
Supposons maintenant que $A_*^*$ soit un complexe\inserted{\unsure{de demi-simpliciens|de demi-simpliciaux|dans la $\frac{1}{2}$ simpliciale}} satisfaisant

% lower, three conditions and qualification to a)
\[
\begin{aligned}
\text{a)}\quad &A^i_n=0\text{ si }i>n
\quad\substack{\text{(il suffit }\deleted{\unsure{A|d}}A^n_n=Z^n_n\text{ i.e.}\\d(A^n_n)=0\ \unsure{\&|!}\\\text{si }n\geq1\text{)}}\\
\text{b)}\quad &H^n(A^*_n)=0\text{ si }n\geq 1\\
\text{c)}\quad &H^{n-1}(A^*_n)=0\text{ si }n\geq 2
\end{aligned}
\]

% foot, conclusion continues
Alors on aura \unsure{qu'un|que un} syst. $K_n:A_n^n\longrightarrow M$ satisfaisant la condition de Stokes

% lower left
% archive sheet number: 66
\scanpage{68}

% top, continuation
est déterminé par la connaissance de

% top, coefficient at zero
\[
\unsure{K_0|k_0|\kappa_0}:A^0_0=A_0\longrightarrow M
\]

% upper, converse
et $K_0$ correspond à un $(K_n)$ \unsure{ssi|sss}

% upper, kernel condition
\[
K_0(\partial_1(\operatorname{Ker}(A^0_1\xrightarrow{d}A^1_1))=0
\]

% upper to middle, universal system
(donc il y a un syst. d'intégrales de Stokes-De Rham ``universel'' à valeurs dans $A_0/\partial_1(\operatorname{Ker}(A_1^0\xrightarrow{d}A_1^1))$ ; le cas le plus intéressant est celui où $\partial_1(\operatorname{Ker}(A_1^0\mathrel{\overset{d}{\unsure{\relbar|\longrightarrow}}}A_1^1))=0$, donc où le choix universel est pour $M=A_0$, -- ce qui sera le cas en particulier si on a

% middle, condition c prime
\[
\text{c')}\quad H^0(A^*_1)\xleftarrow{\sim}\varepsilon^*(A_0^0)\quad(\text{où }\varepsilon:\Delta_1\to\Delta_0\ldots)
\]

% left opposite condition c prime
% slanted
\margin{ou $H^0(A^*_1)=\varepsilon^*(Z^0_0)$ dans le cas où on prend a) sous la forme \unsure{gribouillée|griffonnée}!}

% middle, variant and uniqueness
(ce qui est une variante de c) pour $n=1$).

Sous ces conditions, on voit donc qu'il existe un \uline{unique} hom de complexes de cochaines $\frac{1}{2}$ simpliciaux

% lower middle, map of complexes
\[
A^*_{*}\longrightarrow C^*_{*M}.
\]

% lower middle, degree zero
qui sur $A_0^0=A_0$ coïncide avec

% lower, zero component
\[
K_0:A_0^0\longrightarrow \unsure{C^0_{0M}|\mathbb C^0_{0M}}=M
\]

% lower, nondegenerate cochains
(donc l'identité, à toutes fins pratiques).

On peut noter qu'en fait cet hom est néc. à valeurs dans $C^{*\,!}_{*M}$ (les cochaines non dégénérées), i.e. on a en fait

% foot, exterior power
\[
A^*_\bullet\longrightarrow C^{*\,!}_{*M}\simeq\Lambda^{*+1}\Phi_{*k}\otimes_k M
\]

% foot, last words
en se rappelant que

% bottom edge, intersection of kernels
\[
C^{n\,!}_{*M}=\bigcap_{\sigma:\Delta_n\to\Delta_m,\ m<n}\operatorname{Ker}(C^n_{*M}\xrightarrow{\sigma^*} C^m_{*M})
\]

% lower left
% archive sheet number: 67
\scanpage{69}

% top, connector
donc

% top, Hom and defining condition
\[
\operatorname{Hom}(A^n_*,C^{n\,!}_{*M})=
\{A^n_*\xrightarrow{f^{\unsure{*|n}}}C^n_{*M}\mid \sigma^*\circ f=0\quad\forall\sigma:\Delta_n\to\Delta_m\text{ avec }m<n\}
\]

% upper left, small square
\margin{\diagram{\resizebox{\linewidth}{!}{\begin{tikzcd}[ampersand replacement=\&,column sep=small,row sep=small]
A^n_\bullet\arrow[r,"f^n"]\arrow[d,"\sigma_A^*"'] \& C^n_*\arrow[d,"\sigma_C^*"]\\
A^m_\bullet\arrow[r] \& C^m_*
\end{tikzcd}}}}

% upper, given map
Or \deleted{\unsure{A}} $A^n_\bullet\mathrel{\unsure{\relbar|\longrightarrow}}C^n_\bullet$ est donné par

% upper, map labelled K
\[
A^n_n\xrightarrow{\unsure{K_n|k_n|\kappa_n}}M
\]

% upper, composite
son composé $A^n_\bullet\relbar C^n_\bullet\relbar C^m_*$ est donné par \deleted{la}

% upper middle, M-valued map
\[
A^n_m\longrightarrow M
\]

% upper middle, description
qui n'est autre que le composé

% upper middle, factorisation
\[
A^n_m\xrightarrow{\sigma^*}A^n_n\xrightarrow{K_n}M
\]

% middle, cancellation and new opening
\deleted{et la condition $f(A^n$\unsure{)}} on a donc

% middle, criterion and bound
\[
f^n(A^n_*)\subset C^n_\bullet\quad\Longleftrightarrow\quad
\left|\begin{aligned}
&K_n\circ\sigma^*=0\\
&\forall\sigma:\Delta_n\longrightarrow\Delta_m,\\[-2pt]
&\hspace{2em}\underline{m<n}
\end{aligned}\right.
\]

% middle, degeneracies and reference to a)
(il suffit $=$ de \unsure{l'écrire|la vérifier|le vérifier} quand $\sigma^*$ est une \unsure{op.|opér.} de dégénérescence --) ; or cette condition est tjrs satisfaite si $A_m^n=0$ si $m<n$, ce qu'on a supposé dans a).

% lower, example
\uline{Ex} Soit $S$ un anneau à \unsure{pd-id.|idéal|p.d.} $J$ à p.d., $t\in J$, considérons alors \unsure{l'an$^{\mathrm t}$|cet|l'} Anneau différentiel $\frac{1}{2}$ simplicial à p.d.

% lower, divided-power quotient
\[
A^*_{*S}=(\Gamma^*(\Phi_{*S})\otimes_S\Lambda^*\Psi_{*S})/(T-t.)_{\mathrm{pd}},
\]

% lower, resolution and integrals
c'est une \uline{résolution} de $S$, et $A_j^i=0$ si $i>j$, donc on a un syst. d'intégrales de Stokes-De Rham à valeurs dans $S$, d'où

% bottom, isomorphism and diagonal arrow
\diagram{\begin{tikzcd}[ampersand replacement=\&,column sep=small,row sep=small]
A^*_\bullet\arrow[r,"\sim"] \& \bigl(\Gamma^*(\Phi_{*S})\otimes_{S}\Lambda^*\Psi_{*S}/(T-t.)_{\mathrm{pd}}\bigr)\arrow[dr]\\
\& \& \Lambda^{*+1}_{S}\Phi_{*S}
\end{tikzcd}}

% bottom edge, final remark
dont il faudrait étudier les propriétés multiplicatives\deleted{\unsure{app.}}\ldots

% lower left
% archive sheet number: 68
\scanpage{70}

% top, item 1
\textcircled{1}\quad $k_*\xrightarrow{\varepsilon_*}A_*^*$, $\unsure{K_*|k_*|\kappa_*|K_X}=(K_n)_{n\geq0}\quad K_n:A^n_n\mathrel{\unsure{\relbar|\longrightarrow}} k$, satisf. $\left\{\begin{aligned}&\text{Stokes}\\&K_0\varepsilon_0=\mathrm{id}\end{aligned}\right.$

% upper, item 2
\textcircled{2}\quad Sur \deleted{\unsure{espace|espa}} espace topologique $X$, faisceau différentiel \inserted{de $S$-\unsure{modules|modules alg.}} \deleted{\unsure{$E$|$\varepsilon$}} $\Omega^*$ qui soit une \uline{résolution} de $k_X$ (p.ex. complexe de DR sur $X$ variété $C^\infty$, $k=\mathbb{R}=S$)

% upper middle, item 3, qualifying clauses
\textcircled{3}\quad Un sous-ensemble $\frac{1}{2}$ simplicial $\unsure{S(X)'|\mathbb S(X)'}$ du complexe singulier de $X$, tel que $S(X)'\longrightarrow S(X)$ \deleted{\unsure{soit}} induise un isom sur les groupes d'homologie entière (p.ex. $X$ une variété $C^\infty$ et $S(X)'$ ens. formé des simplexes $C^\infty$) et il en \unsure{soit de même|soit} si on \deleted{\unsure{en}} \unsure{remplace} $X$ par ouvert (d'une base) et qu'on \unsure{prend} de $S(X)$ et $S(X)'$ les simplexes de supp. $\subset U$.

% sloping left margin at item 3
% slanted
\margin{donc, si on regarde le faisceau de coch. sing. $S^*_{X,k}$ sur $X$, et de $=$ $S'_{X,k}$, alors $S'_{X,k}$ comme $S^*_{X,k}$ est une résolution de $k$}

% middle, item 4, overwritten transition and insertions
\textcircled{4}\quad $\forall\sigma\in S(X)'$, et \deleted{\unsure{dont on se donne}} \inserted{hom. \deleted{\unsure{de $k$-alg.}}\inserted{de complexes}} $k$-augm.

% middle, pullback on germs
\[
\sigma^*:\varinjlim_{\substack{U\ \text{ouvert de }X\\ U\supset\operatorname{supp}\sigma}}\Omega^*(U)\longrightarrow A_n^*
\quad(n=\dim\sigma)
\]

% middle, compatibility assumption
On suppose que si $\alpha:\Delta_m\longrightarrow\Delta_n$, $\sigma\in S'(X)_n$, \deleted{$\unsure{\sigma}:\lvert\Delta_n\rvert\longrightarrow X$} on a commutativité dans

% center lower, restriction square
\diagram{\begin{tikzcd}[ampersand replacement=\&,row sep=large]
\displaystyle\varinjlim_{U\supset\operatorname{supp}\sigma}\Omega^*(U)\arrow[r]\arrow[d,"\text{restriction}"] \& A_n^*\arrow[d,"\sigma_A^*"]\\
\displaystyle\varinjlim_{V\supset\operatorname{supp}\alpha^*(\sigma)}\Omega^*(V)\arrow[r] \& A_m^*
\end{tikzcd}}

% left of restriction square
\margin{\uline{NB} $\operatorname{supp}\alpha^*(\sigma)\subset\operatorname{supp}\sigma$}

% lower, bracketed NB
[ \uline{NB} On peut dire que $\sigma\longmapsto\varinjlim_{U\supset\operatorname{supp}\sigma}\Omega^*(U)$ forme\inserted{$\Omega^*_{S(X)'}$} un \deleted{\unsure{faisc.}} objet semi-simplicial (en fait, un complexe de $S$-modules $k$-augmentés) sur $S(X)'=\mathcal{X}$,

% lower left, triangle beside bracketed NB
\margin{\diagram{\resizebox{\linewidth}{!}{\begin{tikzcd}[ampersand replacement=\&,column sep=tiny,row sep=tiny]
\Delta_n\arrow[dr,"\sigma"'] \& \& \Delta_m\arrow[ll,"\alpha"']\arrow[dl,"\alpha^*(\sigma)=\sigma\circ\alpha"]\\
\& S(X)'=\mathcal X \& {}
\end{tikzcd}}}}

% foot, end of bracketed note and final NB
et on se donne un hom de celui-ci dans le complexe analogue $A^\ast_{*\mathcal X}$ sur $\mathcal X=S(X)'$ défini par $A^\ast_{*\mathcal X}$

\uline{NB} finalement \textcircled{1} $(A_*^*,\deleted{\unsure{R|\varepsilon}}k_*,K_*)$ et (4) ne servent que pour pouvoir définir $\unsure{\hookrightarrow|\Leftarrow|\text{un}}\Omega^*\longrightarrow S^{*\,\prime}_X$ comp. \unsure{co-augm.|coaugm.|aux augm.}

% lower left
% archive sheet number: 69
\scanpage{71}

% top center and top right
\mspage{- 140 -}\hfill IX

% upper, definition and first relation
où par définition $({}_nA_K)^f=({}_nA^{\unsure{\mathrm{o}|0}})^f_K$ désigne le r.ème composant $(n=\ell^{r+1})$ du pro-groupe de Barsotti-Tate $T_\ell(A_K)^f$. La première des relations (5.5.8) nous fournit alors un isomorphisme canonique

% upper, (11.6.7)
\[
{}_n\Psi_K\simeq\underline{M}_{nK}\qquad\text{où}\quad\underline{M}_n\overset{\mathrm{dfn}}{=}\underline{M}\otimes_{\mathbb{Z}}(\mathbb{Z}/n\mathbb{Z})\ .\tag*{(11.6.7)}
\]

% upper middle, transition
En fait, on a :

% upper middle, lemma heading and assertion
\uline{Lemme} 11.6.8. \uline{Le schéma en groupes} $ {}_n\Psi$ (11.6.6) \uline{sur} $S$ \uline{est étale sur} $S$, \uline{et l'isomorphisme} (11.6.7) \uline{se prolonge de façon unique en un morphisme}

% middle, open immersion
\[
{}_n\Psi\handwritten{\hookrightarrow}\underline{M}_n\tag*{(11.6.8.1)}
\]

% middle, lemma's last line
\uline{qui est une immersion ouverte}.

% middle and lower, proof of lemma
Pour voir que $ {}_n\Psi$ est étale, il reste à vérifier en vertu de (11.6.7) qu'il l'est en les points de sa fibre spéciale, ou encore que cette dernière est étale, ce qui est trivial, la fibre en question étant $ {}_nA_{\mathrm{o}}/{}_nA_{\mathrm{o}}^{\mathrm{o}}\subset A_{\mathrm{o}}/A_{\mathrm{o}}^{\mathrm{o}}=\Phi_{\mathrm{o}}$. Il s'ensuit que $ {}_n\Psi$ est un schéma normal, et comme $\underline{M}_n$ est fini sur $S$, l'existence d'un prolongement (11.6.8.1) en résulte aussitôt. C'est un morphisme de schémas en groupes étales sur $S$, donc son noyau est un groupe étale sur $S$, et comme $ {}_n\Psi$ donc ce noyau est séparé, et sa fibre générique est le groupe unité par (11.6.7), c'est le groupe unité (3.1.3). Donc (11.6.8.1) est un monomorphisme étale, donc une immersion ouverte (EGA IV 17.9.1), ce qui prouve le lemme.

% lower, new paragraph
Posons maintenant

% lower, definition (11.6.9)
\[
{}_n\Phi=({}_nA)^f/({}_nA^{\mathrm{o}})^f\subset{}_n\Psi\quad,\tag*{(11.6.9)}
\]

% lower, continued sentence
de sorte que $ {}_n\Phi$ est un sous-schéma fini ouvert de $ {}_n\Psi$, puisque $({}_nA)^f$

% bottom strip, pencil numeral
\handwritten{458}

% bottom center, printed number on label
458

% lower left
% archive sheet number: 70
\scanpage{72}

% top across sheet, diagonally cancelled comparison; heavily crossed first pair
\diagram{\deleted{\resizebox{\linewidth}{!}{\begin{tikzcd}[ampersand replacement=\&,column sep=small,row sep=small]
\deleted{\unsure{C^*(X,\Omega^*)}}\arrow[r] \& \deleted{\unsure{\mathbb R\Gamma(X,\Omega^*)}} \& \mathbb R^*\Gamma(X,k)\arrow[l,"\substack{\varepsilon_{\Omega^*}\\\sim}"']\arrow[r,"\substack{\varepsilon_{S_X^*}\\\sim}"]\arrow[d,"\sim"]\arrow[dr,dashed,"\varepsilon_{S_X^{\prime *}}"] \& \mathbb R^*\Gamma(X,S_X^*)\arrow[d,"\sim"] \& \Gamma(X,S^*_{X,\deleted{\unsure{k}}})\arrow[l,"\substack{(S_X^*\ \text{\unsure{étant} par Cartan}\\\text{hom. fin)}\\\sim}"']\arrow[d,"\sim"] \& {}=C^*(S(X)_{\unsure{*|0}},k)\arrow[d,"\sim"]\\
\& \& \mathbb R^*\Gamma(X,\Omega^*)\arrow[r,"\sim"] \& \mathbb R^*\Gamma(X,S_X^{\prime *}) \& \Gamma(X,S_X^{\prime *})\arrow[l,"\sim"'] \& {}=C^*(S(X)'_*,k)\\
\& \& C^*(X,\Omega^*)\arrow[u] \& \& \& {}
\end{tikzcd}}}}

% upper middle, NB introducing the augmentation square
\uline{NB} \unsure{l'hyp.|les hom|la donnée} (4) joint avec $\unsure{K_*|k_*}$ définit un hom

% upper middle right, four-object square embedded in the note
\diagram{\begin{tikzcd}[ampersand replacement=\&,row sep=small]
\Omega^*\arrow[r,"\inserted{\text{par }K}"] \& S_X^{\prime *}\\
k\arrow[u]\arrow[r,dash] \& S_X^*\arrow[u]
\end{tikzcd}}

% upper middle, sentence after square
compatible avec les augmentations, d'où carré comm. \deleted{\unsure{(i) on|(i) ou}} (d'isom dans $D(X,S)$) donnant lieu à un cube commutatif

% center, cube with dashed back edges, Cartan and hypothesis labels
\diagram{\resizebox{\linewidth}{!}{\begin{tikzcd}[ampersand replacement=\&,column sep=small,row sep=large]
C^*(X,\Omega^*)\arrow[rr]\arrow[dd,"\deleted{\unsure{\sim|\wr}}"'] \& \& C^*(X,S_X^{\prime *})\overset{\mathrm{dfn}}{=}C(S'_X,k)\arrow[dd,"\sim"] \& {}\\
\& {[C(X,k)]}\arrow[ul,dashed]\arrow[ul,shift left=1.3ex,"\deleted{\unsure{\leftrightarrow}}" description]\arrow[rr,dashed]\arrow[dd,dashed] \& \& C^*(X,S_X^*)=C(\unsure{S_X|S_x},k)\arrow[ul,"\substack{\text{hyp. sur }S'_X\\\sim}"']\arrow[dd, "{\substack{\sim\ (\text{Cartan, }S_X^*\text{ est}\\\text{``homot.\textsuperscript{t} fin'')}}}"]\\
\mathbb R^*\Gamma(X,\Omega^*)\arrow[rr,"\sim"] \& \& \mathbb R^*\Gamma(X,S_X^{\prime *}) \& {}\\
\& \mathbb R^*\Gamma(X,k)\arrow[ul,"\sim"] \& \& \mathbb R^*\Gamma(X,S_X^*)\arrow[ul,"\sim"]\arrow[ll,"\sim"']
\end{tikzcd}}}

% right margin at cube, written vertically
% margin note written sideways
\margin{\deleted{$\longleftarrow\mathbb R^*\Gamma(X,S^*_{X,k})$}}

% below cube, last four lines
La base du cube est formée d'isom, car transformée par $\mathbb R^*\Gamma$ d'un carré d'iso dans une cat. dérivée ; la face \unsure{devant|dev$^{\mathrm t}$|du haut} est formée d'iso, car on le sait déjà pour \unsure{deux entre eux|deux autres}, donc aussi pour le quatrième qui reste.

% lower left
% archive sheet number: 71
\scanpage{73}

% top center and right
\mspage{- 139 -}\hfill IX

% upper, continuation
définir l'isomorphisme canonique (11.6.2).

% upper, inductive limit
Ecrivant $M_\ell\otimes D_\ell$ (resp. $\check{M}_\ell\otimes D_\ell$) comme limite inductive des $M_n$ (resp. des $\check{M}'_n$), où $n$ parcourt les puissances de $\ell$, il revient au même de définir (11.6.2), où des isomorphismes canoniques

% upper, (11.6.4)
\[
{}_n\Phi_{\unsure{\mathrm{o}|0}}\simeq\operatorname{Ker}(\widetilde{u}_n:M_n\longrightarrow\check{M}'_n)\quad,\quad n=\ell^{r+1}\quad,\tag*{(11.6.4)}
\]

% upper middle, after (11.6.4)
comptatibles avec les morphismes de transition évidents, pour $n$ variable\handwritten{\unsure{!}} (où $\widetilde{u}_n$ désigne maintenant, par abus de notations, l'homomorphisme déduit de $\widetilde{u}_\ell:M_\ell\longrightarrow\check{M}'_\ell$ par réduction mod. $n$).

% middle, canonical filtration
Pour ceci, considérons le schéma en groupes $ {}_nA$ sur $S$, qui est plat et quasi-fini sur $S$ (2.2.1), grâce à l'hypothèse de réduction semi-stable. On a une filtration canonique

% right margin at filtration paragraph
\margin{\handwritten{$\unsure{\lrcorner}$ \unsure{15}}}

% middle, (11.6.5), three right-hand hooks
\[
{}_nA\mathrel{\relbar\joinrel\rhook}({}_nA^{\mathrm{o}})^f\mathrel{\relbar\joinrel\rhook}({}_nA^{\mathrm{o}})^t\mathrel{\relbar\joinrel\rhook}0\quad,\tag*{(11.6.5)}
\]

% middle lower, exponents and completeness
où l'exposant $f$ (resp. l'exposant $t$) désigne la ``partie fixe'' (2.2.3) (resp. la ``partie torique'' (5.1)) du schéma en groupe qu'il affecte ; pour la définition de la partie torique dans 5.1, on notera qu'on avait supposé $S$ complet, ce qui est loisible ici (11.0). (On a d'ailleurs vu dans 6.1 que la construction se généralise au cas $S$ hensélien, si on y tient \ldots) Soit

% lower, (11.6.6)
\[
{}_n\Psi={}_nA/({}_nA^{\mathrm{o}})^f\quad,\tag*{(11.6.6)}
\]

% lower, generic fibre and references
c'est un schéma en groupes plat, séparé et quasi-fini sur $S$ (SGA 3 $\mathrm{VI}_A$ 3.2 et $\mathrm{VI}_B$ 9.2) de fibre générique

% lower, generic-fibre equation
\[
{}_n\Psi_K=({}_nA_K)/({}_nA_K)^f\quad,
\]

% bottom strip, faint numerals
\handwritten{\unsure{457}\qquad\unsure{13}}

% bottom center, printed label
457

% lower right
% archive sheet number: 72
\scanpage{74}

% top, setting
On travaille dans le topos \deleted{\unsure{ss}} $\unsure{\mathrm{Ss}|S_s|\mathcal{S}s},=\widehat{\Delta}$, avec $k$ anneau comm.

% upper left, definition in three lines
\[
\begin{aligned}
\Phi_*:\Phi_{*k}&=i_{0*}(k)\\
&=\Phi_{*\mathbb Z}\otimes_{\mathbb Z}k\\
&=\Phi_{*\mathbb Z}\otimes_{\mathbb Z_*}k_*
\end{aligned}
\]

% upper right, sheaf of modules
\unsure{$\}$|\deleted{$T$}} on le considère comme faisceau de $k$-modules, (pas comme faisceau d'anneaux) muni d'une section $T$

% upper, brace and injection
\[
\left\{\begin{aligned}
T&\in\Gamma(\mathrm{Ss},\Phi_{*k})[=\Gamma(\mathrm{Ss},C^{10}_*)]\quad\text{définissant une \uline{injection}}\\
k_*&\hookrightarrow\Phi_{*k}
\end{aligned}\right.
\]

% upper middle, bigraded algebra
On a la $k$-Algèbre bigraduée \deleted{\unsure{à p.d.}} différentielle à p.d. (diff. de bidegré $-1,+1$ \ldots{} \deleted{\unsure{alg.|(alg.}} alg. alternée pour degré ext.,

% upper middle, tensor product
\[
C^{**}_{*k}=\mathcal{C}^{**}_{*k}=\Gamma^*_{k}\Phi_\bullet\otimes_k\Lambda^*\Phi_*
\]

% left margin alongside algebra
% slanted
\margin{Notations $\overline{\Phi}_*$ pour $\Lambda^1\Phi_*$ et $\bar{x}$ pour \deleted{\unsure{él.}} ``élément'' \deleted{\unsure{de}} $\overline{\Phi}$ correspondant à l'``élément'' \deleted{\unsure{corresp.}} $x$ de $\Phi$}

% upper middle, paragraph continuing through total degree
qui est une résolution de $k$. De façon précise, $C^{00}_{*k}=k_*$, et en degré total donné $n>0$, on trouve une suite exacte

% middle, exact sequence over two lines
\[
\begin{aligned}
0\relbar\Gamma^n\Phi^*&\relbar\Gamma^{n-1}\Phi^*\otimes\Lambda^1\Phi^*
\relbar\Gamma^{n-2}\Phi^*\otimes\Lambda^2\Phi^*\longrightarrow\\
&\ldots\longrightarrow\Phi^*\otimes\Lambda^{n-1}\Phi^*\longrightarrow
\Lambda^n\Phi^*\longrightarrow0\longrightarrow0\ldots
\end{aligned}
\]

% middle, terms and differential of T
dont les termes \deleted{sont} \unsure{tous|sont|donc} flasques.

On considère $\underset{\substack{\Vert\\\bar T}}{dT}\in\Gamma(\mathrm{Ss},\deleted{\unsure{\Phi_*}}C^{01}_*)\simeq\Gamma(\mathrm{Ss},\overline{\Phi}_*)$

% middle, map introduced
et \unsure{cet homomorphisme|l'homomorphisme|cet endomorphisme|\deleted{l'homomorphisme}}

% middle lower, multiplication by barred T
\[
\bar T_\wedge:C^{**}_*\longrightarrow C^{**}_{k}\ ,
\]

% middle lower, marginal qualifications and commutator
on note qu'il est de bidegré $(0,1)$,\inserted{de carré nul, $\operatorname{Im}(\unsure{\widehat T_\wedge|\bar T_\wedge})=\operatorname{Ker}(\bar T_\wedge)$} \inserted{anti}commute à $d$ $(d(\bar T.x)=\underbrace{d\bar T}_{\substack{\Vert\\0}}.x-\bar T dx=-\bar Tdx)$, est injectif sur $C^{*0}$. On pose

% lower, kernels, exterior powers and quotient
\[
\begin{aligned}
\unsure{\mathrm{DR}|\mathbb{D}\mathbb{R}}_{*k}^{pq}=\mathrm{DR}_*^{pq}
 &=\operatorname{Ker}(C^{p,q+1}_*\xrightarrow{\bar T_\wedge}C^{p,q+2}_\bullet)\\
 &=\operatorname{Ker}(\Gamma^p\Phi_\bullet\otimes\Lambda^{q+1}\Phi_\bullet\relbar
 \Gamma^p\Phi_\bullet\otimes\Lambda^{\unsure{q+2|q+1}}\Phi_\bullet)\\
 &\simeq\deleted{\operatorname{Ker}(}\Gamma^p\Phi_\bullet\otimes
 \underbrace{\operatorname{Ker}(\Lambda^{q+1}\Phi_\bullet\relbar\Lambda^{q+2}\Phi_\bullet)}_{\Lambda^q\Psi_*},
 \quad\Psi_\bullet=\Phi_\bullet/k_\bullet\\
 &\simeq\Gamma^p\Phi_\bullet\otimes\Lambda^q\Psi_*\ ,
\end{aligned}
\]

% bottom, induced differential and multiplication
$\mathrm{DR}_*^{**}=(\mathrm{DR}_*^{pq})_{(p,q)\in\mathbb Z\times\mathbb Z}$, avec la bigraduation et la différentielle \inserted{(au signe près)} induites par $C_*^{**}$, \deleted{\unsure{sauf pour}} (mais avec décalage d'une unité pour le degré ext.), et une structure multiplicative \deleted{qui} et à p.d. qui sont caractérisées

% lower left
% archive sheet number: 73
\scanpage{75}

% top, property defining the structures
par la propriété que l'hom épi

% top, epimorphism and inclusion
\[
u=\deleted{\unsure{q}}\bar T_\wedge:\ C^{**}_*\xrightarrow{u}\unsure{\mathrm{DR}|\mathbb D\mathbb R}^{**}_*\mathrel{\overset{i}{\subset}}C^{**}_\bullet
\]

% upper, compatibility
soit compatible à ces structures (en fait, à \uline{ttes} les structures). On aura aussi

% upper, product identities between lines
\[
\deleted{i(ux\,u(y))=}\quad
\left|\begin{aligned}
i(xu(y))&=i(x)y\\
i(u(y)x)&=y\,i(x)
\end{aligned}\right.\qquad i(x)i(y)=0
\]

% upper, exact sequence introduction
Notons qu'on a la suite exacte

% upper middle, short exact sequence, containment beneath target
\[
0\longrightarrow\mathrm{DR}^{p,q}_*\xrightarrow{i}C^{p,q+1}_*
\xrightarrow{(\bar T_\wedge)=u}\underset{\substack{\cap\\C^{p,q+2}_*}}{\mathrm{DR}^{p,q+1}_*}\longrightarrow0
\]

% upper middle, flabby kernel and sections
et que le noyau $\mathrm{DR}^{\unsure{p,q|q,p}}_*$ est flasque si $p\deleted{\unsure{=}}>0$ -- donc si $X_*\in\operatorname{Ob}\mathrm{Ss}$, $\operatorname{Hom}(X_*,-)$ dans cette suite exacte donne encore une suite exacte (i.e. $\deleted{\unsure{C^{p,q+1}}}C^{p,i}(X_*)\xrightarrow{\bar T}\mathrm{DR}^{p,i}(X_*)$ est surjectif).

% middle, passage to Hom
Passons donc à $\operatorname{Hom}(X_*,-)$.

% middle, definition of C
\[
\mathcal C^{**}=C^{**}_*(X_*)=C^{**}(X_*,k)
\]

% middle, algebra and T
est une $k$-Algèbre bigraduée différentielle à p.d., avec élément marqué $T\in C^{1,0}$

% middle lower, (A)
\[
\left\{\begin{aligned}
&C^{00}\inserted{=H^{00}}\simeq k^{\pi_0(X_\bullet)}\deleted{\unsure{X_*}}\quad(=k\text{ si }X_\bullet\text{ connexe}),\\
&H^{pq}(\mathcal C^{**})=0\quad\text{si }pq\ne0,\quad\relbar
\end{aligned}\right.\tag*{(A)}
\]

% lower, definition preceding (B)
Posant $\bar T=dT\in C^{0,1}$, on a de plus

% lower, (B)
\[
\left\{\begin{aligned}
&0\longrightarrow C^{p,0}\xrightarrow{\bar T_\wedge}C^{p,\deleted{\unsure{q+}}1}
 \xrightarrow{\bar T_\wedge}C^{p,2}\longrightarrow\ldots
 \longrightarrow C^{p,q}\mathrel{\overset{\bar T_\wedge}{\unsure{\relbar|\longrightarrow}}}C^{p,q+1}\longrightarrow\\
&\text{est une suite exacte si }p>0
\end{aligned}\right.\tag*{(B)}
\]

% lower, exceptional column
tandis que, pour $p=0$, comme

% lower, exterior-power sequence and identification below first term
\[
\deleted{\unsure{\Phi}}0\relbar\underset{\substack{\Vert\\k_\bullet}}{\Lambda^0\Phi_\bullet}
\xrightarrow{\bar T_\wedge}\Lambda^1\Phi\mathrel{\overset{\bar T_\wedge}{\relbar}}\Lambda^2\Phi_\bullet\longrightarrow
\]

% lower, flabby resolution
est une suite exacte fournissant une résolution flasque de $k_*$, on trouve

% bottom, cohomology computed in two cases
\[
\begin{aligned}
&H^{\deleted{\unsure{*}}n}(0\relbar C^{0,0}\xrightarrow{\bar T_{\unsure{\wedge|\bullet}}}C^{0,1}\longrightarrow\ldots
\longrightarrow C^{0,q}\xrightarrow{\bar T_\wedge}C^{0,q+1}\longrightarrow\ldots)\\
&\simeq\begin{cases}0&\text{si }n\leq1\\ H^{n-1}(X_*,k)&\text{si }n\geq2\end{cases}
\end{aligned}
\]

% lower left
% archive sheet number: 74
\scanpage{76}

% top, introduction
Posant

% top, DR and its graded family
\[
\unsure{\mathrm{DR}|\mathbb D\mathbb R}^{pq}=\mathrm{DR}^{pq}(X_*)=\operatorname{Ker}(C^{p,q+1}\xrightarrow{\bar T_{\unsure{\wedge|\bullet}}}C^{p,q+2})\ ,\qquad
\mathrm{DR}^{**}=(\mathrm{DR}^{pq})_{p,q\in\unsure{\mathbb Z|\mathbb Z\times\mathbb Z}}
\]

% upper, transition
on aura donc

% upper, two stages of the map and inclusions
\[
\underset{\text{épi si }p\ne0}{\smash[b]{\mathstrut C^{p,q}\xrightarrow{\bar T_.\,=u}\underset{\substack{\cap\\C^{p,q+1}}}{\mathrm{DR}^{pq}}}}\qquad
\text{d'où}\quad C^{**}\xrightarrow{u}\underset{\substack{\cap\,i\\ C_*^{**}}}{\mathrm{DR}^{**}}\quad\deleted{\unsure{u}}
\]

% upper left, boxed marginal identities
\margin{$\begin{aligned}\operatorname{Im}i&=\operatorname{Ker}u\\i\circ u&=\bar T_\wedge(\,.\,)\end{aligned}$}

% upper, (C)
\[
i(x).i(y)=0\tag*{(C)}
\]

% upper, parenthesis and degree restriction
(il est évident que $i(\mathrm{DR})\underline{i(u\,\mathrm{DR})=0}$ donc $i(\mathrm{DR}^{p,q})\,i(\mathrm{DR}^{p',q'})=0$ si $p$ ou $p'\ne0$)

% left, oblique boxed marginal reformulation
% slanted
\margin{donc (C) équiv$^{\mathrm t}$ à $\boxed{i(\mathrm{DR}^{0p}).i\mathrm{DR}_*^{0q}\mathrel{\unsure{=|\ne}}0}$}

% upper middle, algebra structure
\deleted{et} on aura sur $\mathrm{DR}^{**}$ une structure d'algèbre graduée différentielle à p.d. \inserted{canonique, avec axiomes usuels (alternée etc)} \deleted{\unsure{telle que}}

% upper middle, map u
\[
u:C^{**}\mathrel{\unsure{\relbar|\longrightarrow}}\unsure{D|\mathrm{DR}}^{**}
\]

% middle, compatibility and qualifications
soit compatible avec ttes ces structures \inserted{et $\left\{\begin{aligned}i(xuy)&=i(x)y\\i(uyx)&=y\,i(x)\end{aligned}\right.$}. Ces structures sur $\mathrm{DR}^{**}$ sont presque déterminées par ces propriétés, plus celle que l'inclusion $\mathrm{DR}^{**}\subset C^{**}$ est de bidegré $0,1$ et \inserted{anti}commute à la différentielle -- à la détermination près des produits $xy$ \deleted{\unsure{et $x$}} et des puissances divisées $x^{(i)}$ $(i\geq2)$ quand $x,y$ \deleted{sont} de degré complémentaire $0$, (auquel cas on ne peut pas néc. les écrire sous la forme $u(x')$\inserted{\deleted{\unsure{et}} ou $u(y')$}, \unsure{où|tiré|tels} dans $C^{**}$).

% middle lower, bracketed NB and crowded revisions
[ \uline{NB} Il est évident que sur \deleted{\unsure{$\mathrm{DR}$}} on peut définir \inserted{de façon unique} une multiplication \inserted{bilinéaire sur $\mathrm{DR}\times\mathrm{DR}$,\inserted{$(\operatorname{Im}u)\times\mathrm{DR}$ et $\mathrm{DR}\times\operatorname{Im}u$}\deleted{\unsure{$\in k$}}} \deleted{\unsure{et}} satisfaisant

% lower middle, three identities
\[
u(x.y)=u(x)u(y),\qquad i(xu(y))=i(x)y,\quad i(u(y)x)=y\,i(x)
\]

% lower middle, implication and applying i
(\deleted{\unsure{car}} la première relation équivaut, en appliquant $i$, à $\bar T_.(xy)=i(ux.uy)$, i.e. en appliq. la deuxième $\bar T.(xy)=(\bar T.x)y$

% lower, two maps and existence
donc il suffit de vérifier l'existence et l'unicité de $\operatorname{Im}u\times\mathrm{DR}\relbar\mathrm{DR}$ \inserted{et $\mathrm{DR}\times\operatorname{Im}u\relbar\mathrm{DR}$} satisfaisant les rel. 2 et 3. Or pour la $2^{\mathrm{ème}}$, il suffit de noter que

% lower, a)
a) $i(x)y$\inserted{$=xy$} $\in\deleted{\operatorname{Ker}\bar T}\mathrm{DR}=\operatorname{Ker}(\bar T_\wedge)$ car $\bar T(xy)=(\bar T x)y=0$ car $\bar T x=0$ puisque $x\in\mathrm{DR}$

% foot, b)
b) $i(x)y$\inserted{$=xy$} ne dépend que de $x\in\mathrm{DR}$ et de $uy\in\mathrm{DR}$, car si $\deleted{\unsure{u}}u\unsure{y)|(y)}=0$ i.e. $\bar Ty=0$, i.e. $y\in\mathrm{DR}$, on a $xy=0$ -- Idem pour 3.

% lower left
% archive sheet number: 75
\scanpage{77}

% top, operations on the image
De $=$, on a sur $u(C^{**+})\subset\unsure{\mathrm{DR}|\mathbb D\mathbb R}^{**+}$ des op. ``puiss. div'' \inserted{uniques} données par la condition

% top, compatibility with u
\[
u(x)^{(i)}=u(x^{(i)})
\]

% upper, change of representative and interlinear version
car il suffit de prouver que si $x'$ est tel que $ux=ux'$, i.e. $x'=x+y$, $\bar T.y=0$, alors $u(x'^{(i)})=u(x^{(i)})$, \inserted{i.e. $\bar T_\wedge x'^{(i)}=\bar T x^{(i)}$} or

% upper, expansion
\[
x'^{(i)}=x^{(i)}+\deleted{\unsure{u}}x^{(i-1)}y+x^{(i-2)}y^{(2)}\mathrel{\relbar}+y^{(i)}
\]

% upper middle, reduction to D
il suffit de prouver que $\bar T y^{(i)}=0$ pour $i\geq1$, ce qui résulte du fait \deleted{\unsure{suivant}} \deleted{\unsure{C}}

% upper middle, D, with struck wording
\textcircled{D}\quad \inserted{sur} $\underset{\substack{\Vert\\\operatorname{Ker}(\bar T_\wedge(.))}}{\mathrm{DR}^{**\deleted{\unsure{+}}}}\subset C^{**}$ \deleted{est stable par} les puiss. div. (de $C^{**}$) $y^{(i)}$ $(i\geq2)$ sont nulles

% upper middle, cancelled restatement of D
\deleted{i.e. si $y\in C^{**+}$, $\bar T.y=0\ \Longrightarrow\ \bar T.y^{(i)}=0\quad\forall i\geq1$,}

% middle, local calculation
(ceci, \unsure{à son tour|son tour|se démontre} résulte de la construction pour $\mathrm{DR}^{**}_\bullet\subset C^{**}_*$, où loc. $y$ est de la forme $\bar Tz$, donc $\unsure{\bar T|\deleted{\bar T}}y^{(i)}=\bar T^i z^{(i)}$, qui est donc nul si $\unsure{i|\partial}\geq2$.)

% middle, heavily crossed display
\deleted{$(\unsure{E|F})\quad\unsure{T^{(n)}|i(}\bar T(\unsure{z}\,\mathrm{DR}^{n,0}\subset C^{n,1})\unsure{\subset u|=u}$}

% middle lower, additional data
Ainsi, les seules données supplémentaires sur $C^{**}$ pour \deleted{définir}\inserted{décrire} les structures habituelles sur $\mathrm{DR}^{**}$ sont

% lower, E, two operations, revised symbols and exponents
\[
\substack{\text{(E)}\\\text{(\uline{données})}}\quad
\left\{\begin{aligned}
&\text{``multiplication''}\quad \deleted{\unsure{Z|\mathrm{DR}}}\mathrm{DR}^{0,p\deleted{\unsure{+1|+}}}
 \times\deleted{Z}\mathrm{DR}^{0,q\deleted{+1}}\longrightarrow\deleted{Z}\mathrm{DR}^{0,p+q\deleted{+1}}\quad(p,q\geq1)\\
&\text{``puiss. divisées''}\quad \deleted{Z}\mathrm{DR}^{0,p}\xrightarrow{\gamma^i}\deleted{Z}\mathrm{DR}^{0,pi\deleted{\unsure{+1}}}\quad(i\geq2)
\end{aligned}\right.
\]

% lower, boxed and cancelled NB
\deleted{(\uline{NB} \unsure{dans} $Z^{0,1}\times Z^{0,q+1}\relbar$ pour $p$ ou $q=1$, on a, \deleted{\unsure{$Z^{0,p+1}$}} puisque $Z^{0,1}$ est engendré par}

% foot, axioms
plus les axiomes disant qu'il faut pour assurer que sur $\mathrm{DR}^{**}$, on a les propriétés habituelles pour les opérations.

% lower left
% archive sheet number: 76
\scanpage{78}

% heading
\uline{Pb} trouver des splittings \uline{canoniques} des suites exactes

% upper half, vertical sequence and proposed section
\diagram{\begin{tikzcd}[ampersand replacement=\&,row sep=small]
0\arrow[d] \& \\
\Gamma^p\Phi_*\otimes\Lambda^q\Psi_*\arrow[r,dashed,dash]\arrow[d,dashed,bend left=65,"{?}"] \& \Lambda^{q+1}\Phi_*\\
\Gamma^p\Phi_*\otimes\Lambda^q\bar\Phi_*\arrow[u] \& \\
\Gamma^p\Phi_*\otimes\Lambda^{q-1}\Psi_*\arrow[u]\arrow[r,dashed] \& \Lambda^q\Phi_*\\
0\arrow[u] \& \text{pour }p\geq1\ .
\end{tikzcd}}

% below the diagram, with rising caret after projectif
\uline{NB} Comme $\Gamma^p\Phi_*\otimes\Lambda^q\Psi_*$ est un $k$-Module projectif \inserted{$(p\geq1)$}, il y a un splitting. On veut un choix simple. Noter que le cas $q=0$ est trivial (car $\Lambda^{q-1}\Psi_*=0$)

% middle heading
\uline{Cas premier} $p=1$, $q=1$

% middle, special sequence
La suite exacte est

% middle, sequence with identification underneath
\[
0\relbar\underset{\substack{\Vert\\\Phi_*}}{\Phi_*\otimes k}\longrightarrow\Phi_*\otimes\bar\Phi_*\relbar\Phi_*\otimes\Psi_*\relbar0
\]

% below the sequence
\unsure{NB|On} note que

% lower middle, Hom comparison with annotations underneath
\[
\begin{aligned}
\operatorname{Hom}_{k\text{-Mod}}(\Phi_*\otimes\Psi_*,\underset{\substack{\Vert\\i_{0*}(k)}}{\Phi_*})
 &\simeq\operatorname{Hom}_{\unsure{k\mathrm{mod}|k\text{-mod}|k\text{-Mod}}}(\underset{\Phi_0\otimes\underset{\substack{\Vert\\0}}{\Psi_0}=0}{(\underline{\Phi_*\otimes\Psi_*})_0},k)\\
 &=0,
\end{aligned}
\]

% lower, uniqueness
donc splitting \uline{unique}

% lower, direct sum
\[
\Phi_*\otimes\bar\Phi_*\simeq(\Phi_*\otimes\Psi_*)\oplus\Phi_*
\]

% last line, opening with cancellation
Soit \deleted{\unsure{un|s}}

% last line, proposed s
\[
s:\Phi_*\otimes\Psi_*\relbar\Phi_*\otimes\Phi_*
\]

% lower left
% archive sheet number: 77
